Neural-symbolic learning is often motivated by systematic generalization: a model should learn perceptual predicates from data while using symbolic structure to compose them beyond the observed training labels. We test this claim in a controlled \addmnist setting where training excludes all pairs in which both digits are high, i.e., both in $\{5,\ldots,9\}$. This restriction removes sums 14--18 from the supervised training labels while preserving the symbolic rule for addition. We compare a direct pair-image CNN against three neural-symbolic models that learn digit predicates and compute pair sums by exact marginalization over the addition table, including two anchor-based variants. With 5,000 weak pair labels, the plain symbolic likelihood model reaches 90.0\% IID accuracy and 82.5\% OOD high-plus-high accuracy, while the direct CNN reaches 31.8\% IID accuracy and 0.3\% OOD accuracy. On the truly unseen OOD sums 14--18, the symbolic likelihood model reaches 80.6\% accuracy, whereas the direct CNN reaches 0.0\%. The proposed anchor-posterior variant improves strongly over the direct baseline but does not beat the simpler symbolic likelihood. These results show that exact symbolic marginalization can extrapolate to labels never observed during training, but only after the neural digit predicates become well grounded.
